\documentclass[10pt]{article}
\usepackage[margin=0.7in]{geometry}
\usepackage{amsmath}
\usepackage{booktabs}
\usepackage{array}
\usepackage{hyperref}
\usepackage{float}

\title{GridFM as an AC power-flow surrogate\\
\large 31 PPC grids, one model per grid}
\author{PIGNN-Attn-LS-PPC experimental summary}
\date{11 August 2026}

\begin{document}
\maketitle

\section{Scope}

One AC power-flow surrogate trained per grid across the Phase-0 backbone corpus:
given the bus injections, predict the voltage state. This is the power-flow task
(\texttt{-{}-task pf}), not OPF --- no decision space, no dispatch, no cost.

\paragraph{Which GridFM this is.} The model here is
\texttt{GridFMHeteroSurrogate}, a local re-implementation in GridFM's
architectural style, not IBM's released \texttt{gridfm\_graphkit} model. Nothing
imports the released model and no checkpoint is loaded, so every run is from
scratch. The released \texttt{GNS\_heterogeneous} adds per-layer physics
feedback and a \texttt{PhysicsDecoderPF} that recovers the quantities the
balance equations determine; a matched run of that model is in progress and is
not reported here. \textbf{Do not read these numbers as ``GridFM's''
performance} --- they are this mirror's.

\paragraph{Data.} The Phase-0 backbone corpus: the pandapower-solved scenario
set tagged \texttt{\_ppNR\_}, plus LVN\_heo1, which carries \texttt{\_cNR\_}
because pandapower cannot solve that grid and it was produced with the custom
Newton oracle. 31 grids, 4 to 9241 buses, mostly 36{,}000 scenarios each, split
$1/3$ train, $1/3$ valid, $1/3$ test.

\paragraph{Setup.} 40 epochs, AdamW at LR $5\times10^{-4}$, hidden 48, 12
layers, 8 heads (the released config's capacity), loss
$\mathrm{MSE}(V) + 10^{-2}\,\mathcal{L}_{\rm phys}$ with a log-cosh physics
form. Per-unit on a 100~MVA base, complex128 throughout. Run on helma H100s,
one job per grid.

\section{Two changes the corpus forced}

\paragraph{Row groups.} Most \texttt{*\_NR\_branchrows\_directSI.parquet} files
were written with 3000--6000 rows per group. Training shuffles, so a batch of a
few samples decodes a whole group; the same effect measured 3191~ms/batch
against 176~ms/batch on the OPF files, an $18\times$ penalty. 28 files were
re-laid-out to 20 rows per group before training. SimBench, ENTSO-E and
LVN\_heo1 were already at 20.

\paragraph{A sparse Y-bus.} The collate built a \emph{dense} block-diagonal
Y-bus, costing $(B\!\cdot\!N)^2\times16$ bytes --- quadratic in batch times
buses. case2848rte at batch 32 asked for 124~GB and every large grid was
OOM-killed. The matrix is 99.98\% zeros, so it is now assembled as sparse COO:

\begin{table}[H]
\centering
\small
\begin{tabular}{lrr}
\toprule
case2848rte, batch 8 & dense & sparse \\
\midrule
Y-bus footprint & 7921 MB & \textbf{1.8 MB} \\
nonzeros & --- & 77{,}856 (0.015\%) \\
\bottomrule
\end{tabular}
\caption{Verified numerically identical, not merely close:
\texttt{to\_dense(sparse)} equals the dense matrix exactly, and the implied
injection $S = V\odot\overline{YV}$ agrees to $1.1\times10^{-11}$ absolute
($5.2\times10^{-13}$ relative), i.e. complex128 rounding. Gradients flow
unchanged.}
\end{table}

With the Y-bus out of the way the binding constraint became the model's own
activation memory. Measured directly on case9241pegase --- 4.3~GB at batch 1
rising linearly to 67.0~GB at batch 16, OOM at batch 24 --- which works out to
$\sim$479~KB per node and sets the per-grid batch (12 for case9241pegase, 64 for
everything up to case1354pegase).

\section{Results: voltage}

\begin{table}[H]
\centering
\small
\begin{tabular}{lrrrrr}
\toprule
Grid & Buses & RMSE & $|V|$ RMSE & $\theta$ [deg] & batch \\
\midrule
SimBench & 94 & 1.203e-03 & 1.05e-03 & 0.034 & 64 \\
case4gs & 4 & 2.367e-03 & 1.69e-03 & 0.095 & 64 \\
case5 & 5 & 3.666e-03 & 2.61e-03 & 0.148 & 64 \\
case6ww & 6 & 4.948e-03 & 1.56e-03 & 0.269 & 64 \\
case30 & 30 & 5.545e-03 & 2.41e-03 & 0.286 & 64 \\
case14 & 14 & 5.829e-03 & 2.50e-03 & 0.302 & 64 \\
case\_ieee30 & 30 & 7.684e-03 & 4.53e-03 & 0.356 & 64 \\
case9 & 9 & 8.907e-03 & 3.83e-03 & 0.461 & 64 \\
case24\_ieee\_rts & 24 & 9.392e-03 & 2.06e-03 & 0.525 & 64 \\
case\_illinois200 & 200 & 1.074e-02 & 6.62e-03 & 0.484 & 64 \\
case33bw & 33 & 1.138e-02 & 1.10e-02 & 0.170 & 64 \\
case118 & 118 & 1.144e-02 & 2.94e-03 & 0.633 & 64 \\
case57 & 57 & 1.177e-02 & 8.65e-03 & 0.457 & 64 \\
case39 & 39 & 1.511e-02 & 6.96e-03 & 0.768 & 64 \\
case6495rte & 6495 & 2.603e-02 & 1.77e-02 & 1.094 & 18 \\
LVN\_heo1 & 722 & 2.756e-02 & 2.74e-02 & 0.187 & 64 \\
GBnetwork & 2224 & 3.156e-02 & 2.26e-02 & 1.263 & 53 \\
GBreducednetwork & 29 & 3.188e-02 & 2.29e-02 & 1.271 & 64 \\
case89pegase & 89 & 3.374e-02 & 2.73e-02 & 1.139 & 64 \\
iceland & 189 & 3.606e-02 & 1.71e-02 & 1.820 & 64 \\
case2848rte & 2848 & 3.720e-02 & 2.61e-02 & 1.519 & 42 \\
case6515rte & 6515 & 4.151e-02 & 2.25e-02 & 1.997 & 18 \\
case300 & 300 & 4.835e-02 & 1.15e-02 & 2.691 & 64 \\
case1354pegase & 1354 & 5.199e-02 & 2.41e-02 & 2.638 & 64 \\
case6470rte & 6470 & 6.105e-02 & 2.17e-02 & 3.270 & 18 \\
case1888rte & 1888 & 6.878e-02 & 3.06e-02 & 3.529 & 63 \\
case2869pegase & 2869 & 7.219e-02 & 2.52e-02 & 3.876 & 41 \\
case3120sp & 3120 & 1.127e-01 & 2.84e-02 & 6.246 & 38 \\
case9241pegase & 9241 & 1.227e-01 & 2.76e-02 & 6.849 & 12 \\
case145 & 145 & 2.060e-01 & 2.76e-02 & 11.695 & 64 \\
\bottomrule
\end{tabular}
\caption{Final test-set voltage error, sorted by RMSE. ENTSO\_E\_RealGridTest
had not finished at the time of writing and is omitted. RMSE combines magnitude
and angle; the two are also given separately because they behave very
differently.}
\end{table}

\section{Results: power-balance residuals}

Voltage error says how close the predicted state is to the reference one. It
does not say whether that state satisfies the network equations, which is what
an operator would care about. The residual
$\Delta S = S_{\rm implied} - S_{\rm set}$ is computed from
$S = V\odot\overline{YV}$ and masked by bus type: active power is not scored at
the slack and reactive power is not scored at slack or PV buses, because there
the injection is a free variable rather than a known input.

\begin{table}[H]
\centering
\small
\begin{tabular}{lrrrrrr}
\toprule
Grid & $\Delta P^{\infty}$ & $\Delta Q^{\infty}$ & $\overline{|\Delta P|}$ & $\overline{|\Delta Q|}$ & p95$|\Delta P|$ & p95$|\Delta Q|$ \\
\midrule
SimBench & 0.0267 & 0.0396 & 0.00336 & 0.00665 & 0.0119 & 0.0241 \\
case33bw & 0.0509 & 0.0592 & 0.00893 & 0.00854 & 0.031 & 0.0276 \\
case30 & 0.0935 & 0.0874 & 0.0203 & 0.0165 & 0.0671 & 0.0512 \\
case6ww & 0.039 & 0.0233 & 0.0205 & 0.0141 & 0.0522 & 0.037 \\
case\_ieee30 & 0.0848 & 0.0742 & 0.0225 & 0.0193 & 0.066 & 0.057 \\
LVN\_heo1 & 3.33 & 23.4 & 0.0246 & 0.125 & 0.0442 & 0.321 \\
case14 & 0.0838 & 0.0493 & 0.0255 & 0.0197 & 0.0791 & 0.0531 \\
case4gs & 0.0431 & 0.0466 & 0.0261 & 0.0344 & 0.064 & 0.0814 \\
case9 & 0.0901 & 0.0818 & 0.04 & 0.0364 & 0.102 & 0.0988 \\
case57 & 0.221 & 0.178 & 0.0434 & 0.0364 & 0.143 & 0.117 \\
case5 & 0.0943 & 0.0568 & 0.0516 & 0.0568 & 0.13 & 0.133 \\
case\_illinois200 & 0.517 & 0.284 & 0.062 & 0.0463 & 0.217 & 0.137 \\
case118 & 0.553 & 0.365 & 0.072 & 0.0534 & 0.232 & 0.162 \\
case24\_ieee\_rts & 0.275 & 0.16 & 0.0806 & 0.0517 & 0.231 & 0.149 \\
case39 & 0.565 & 0.496 & 0.16 & 0.143 & 0.425 & 0.395 \\
case300 & 5.52 & 4.85 & 0.297 & 0.217 & 0.995 & 0.706 \\
case6470rte & 245 & 142 & 0.531 & 0.305 & 1.56 & 0.875 \\
case6515rte & 248 & 141 & 0.552 & 0.34 & 1.51 & 0.956 \\
case6495rte & 340 & 793 & 0.593 & 0.379 & 1.62 & 0.671 \\
case1354pegase & 49.4 & 18.1 & 0.908 & 0.467 & 3.26 & 1.58 \\
case3120sp & 431 & 79.1 & 1.14 & 0.368 & 2.27 & 0.914 \\
case2848rte & 187 & 106 & 1.19 & 0.612 & 3.75 & 1.65 \\
case2869pegase & 105 & 17.7 & 1.23 & 0.395 & 4.08 & 1.29 \\
case9241pegase & 206 & 53.5 & 1.68 & 0.493 & 5.23 & 1.6 \\
case89pegase & 24.6 & 7.71 & 2.15 & 0.845 & 10.1 & 3.3 \\
case1888rte & 264 & 71 & 2.46 & 0.744 & 4.91 & 2.24 \\
case145 & 48.1 & 8.88 & 4.35 & 1.12 & 16.9 & 4.16 \\
GBnetwork & 1.3e+03 & 1.38e+03 & 4.39 & 1.21 & 3.98 & 1.49 \\
GBreducednetwork & 47.2 & 27.4 & 5.08 & 12.4 & 26.5 & 35 \\
iceland & 746 & 97.8 & 8.18 & 0.795 & 0.899 & 0.53 \\
\bottomrule
\end{tabular}
\caption{All quantities in pu on a 100~MVA base, so $10^{-2}$~pu $=$ 1~MW/MVAr.
Sorted by $\overline{|\Delta P|}$.}
\end{table}

\section{Findings}

\paragraph{The residual ranking is not the voltage ranking.} Spearman
correlation between RMSE and $\overline{|\Delta P|}$ is 0.78 over the 30 finished
grids --- strong but far from the identity a single-metric summary would imply.
The disagreements are large and systematic:

\begin{table}[H]
\centering
\small
\begin{tabular}{lrrr}
\toprule
Grid & RMSE rank & residual rank & shift \\
\midrule
GBnetwork        & 17 & 28 & $-11$ \\
GBreducednetwork & 18 & 29 & $-11$ \\
iceland          & 20 & 30 & $-10$ \\
LVN\_heo1        & 16 &  6 & $+10$ \\
case33bw         & 11 &  2 & $+9$ \\
case5            &  3 & 11 & $-8$ \\
\bottomrule
\end{tabular}
\caption{Grids whose rank moves by 8 or more between the two metrics. A model
can track the reference voltages closely and still leave a large injection
mismatch, and the reverse.}
\end{table}

The practical consequence: \textbf{reporting voltage RMSE alone would have rated
iceland (rank 20 of 30) as mid-pack when it has the worst mean active-power
residual in the corpus}, and would have understated LVN\_heo1, which is 16th on
voltage but 6th on residual.

\paragraph{The residual distribution is extremely heavy-tailed.} On two grids
the mean exceeds the 95th percentile --- iceland at
$\overline{|\Delta P|}=8.18$ against p95 $=0.899$, and GBnetwork at 4.39 against
3.98. That ordering is only possible if a small number of buses carry residuals
orders of magnitude above the rest. Reporting a mean alone would hide this
entirely, and reporting $\Delta P^{\infty}$ alone would suggest the whole grid is
bad when it is a handful of buses. Both are needed.

The worst $\Delta P^{\infty}$ values are physically implausible as
\emph{modelling} error: GBnetwork reaches 1300~pu, i.e.\ 130~GW at a single bus.
These are better read as a signal that something degenerate is happening at
specific buses --- most likely near-zero-impedance branches or isolated
components --- than as a smooth accuracy figure.

\paragraph{Absolute error is dominated by the angle.} Across the corpus
$|V|$ RMSE spans $1.0\times10^{-3}$ to $3.1\times10^{-2}$, a factor of 31, while
$\theta$ spans $0.034^\circ$ to $11.7^\circ$, a factor of 344, so the ranking is
set almost entirely by the angle. It does \emph{not} follow that magnitude is
the easier channel: Section~\ref{sec:diag} shows the magnitude is the one that
collapses toward its mean on half the corpus, which a small absolute error
conceals.

\paragraph{Bus count does not predict difficulty; grid family does.}
GBreducednetwork has 29 buses and a worse RMSE ($3.19\times10^{-2}$) than
case118 with 118 buses ($1.14\times10^{-2}$); case6495rte with 6495 buses beats
case300 with 300. What does track is provenance: the IEEE-standard cases occupy
the top of the table and the pegase / rte / sp / GB families the bottom,
independent of size. Whatever makes those grids hard --- ill-conditioning,
topology, voltage-level spread --- is not captured by bus count.


\section{MAE and the collapse diagnostic}
\label{sec:diag}

The tables above report RMSE. Two papers argue that is not enough, and both
turn out to matter here.

\paragraph{Why MAE alone is also not enough.} GridSFM's white paper makes the
point against its own headline metric: bus voltage magnitude in pu spans only
about 0.95--1.10, so \emph{a flat predictor of $V = 1.0$ at every bus already
posts a small MAE} while carrying no information about the voltage profile.
MAPE has the mirror problem, being dominated by the small-denominator tail.
Neither says whether a model follows the reference's variation or collapses
toward its mean. Their remedy, adopted here, is a per-channel pooled regression
of predicted against reference values: slope~1 means the model tracks the
reference spread one-for-one, and slope well below~1 means collapse. We regress
the angle on $\sin\theta$ and $\cos\theta$ rather than on $\theta$, because a
wrapped angle is circular and a raw fit is corrupted by any pair straddling the
branch cut.

This is not a hypothetical failure mode for this project. A GridSFM run on
OPFData case118 was found to have collapsed its magnitude head only because two
independently initialised runs reported a $|V|$ RMSE identical to five
significant figures. A slope near zero would have shown it directly, from one
run.

\begin{table}[H]
\centering
\small
\begin{tabular}{lrrrrrrrr}
\toprule
 & \multicolumn{2}{c}{$|V|$ [pu]} & \multicolumn{2}{c}{$\theta$ [deg]} & \multicolumn{2}{c}{regression $|V|$} & \multicolumn{2}{c}{regression $\sin\theta$} \\
\cmidrule(lr){2-3}\cmidrule(lr){4-5}\cmidrule(lr){6-7}\cmidrule(lr){8-9}
Grid & MAE & RMSE & MAE & RMSE & slope & $R^2$ & slope & $R^2$ \\
\midrule
case3120sp & 2.274e-02 & 2.843e-02 & 4.143 & 6.246 & \textbf{0.208} & 0.2923 & 0.964 & 0.9276 \\
LVN\_heo1 & 2.052e-02 & 2.737e-02 & 0.121 & 0.187 & \textbf{0.213} & 0.2251 & 0.922 & 0.9956 \\
case1354pegase & 1.887e-02 & 2.414e-02 & 1.985 & 2.638 & \textbf{0.326} & 0.4159 & 0.930 & 0.9474 \\
case2869pegase & 1.975e-02 & 2.521e-02 & 2.884 & 3.876 & \textbf{0.330} & 0.4204 & 0.955 & 0.9860 \\
case89pegase & 2.203e-02 & 2.725e-02 & 0.874 & 1.139 & \textbf{0.381} & 0.4700 & 0.948 & 0.9810 \\
case1888rte & 2.443e-02 & 3.060e-02 & 2.583 & 3.529 & \textbf{0.443} & 0.5587 & 0.937 & 0.9639 \\
case2848rte & 2.012e-02 & 2.609e-02 & 1.196 & 1.519 & \textbf{0.621} & 0.7287 & 0.946 & 0.9736 \\
case6470rte & 1.659e-02 & 2.166e-02 & 2.315 & 3.270 & \textbf{0.738} & 0.7612 & 0.996 & 0.9968 \\
GBnetwork & 1.572e-02 & 2.259e-02 & 0.942 & 1.263 & \textbf{0.739} & 0.8252 & 0.990 & 0.9975 \\
case145 & 1.989e-02 & 2.761e-02 & 8.605 & 11.695 & \textbf{0.780} & 0.7865 & 0.992 & 0.9576 \\
case33bw & 7.655e-03 & 1.099e-02 & 0.134 & 0.170 & \textbf{0.872} & 0.8960 & 0.776 & 0.7827 \\
case39 & 5.107e-03 & 6.957e-03 & 0.587 & 0.768 & \textbf{0.889} & 0.9475 & 0.981 & 0.9910 \\
iceland & 1.316e-02 & 1.709e-02 & 1.339 & 1.820 & \textbf{0.897} & 0.9064 & 0.988 & 0.9812 \\
case\_illinois200 & 4.745e-03 & 6.625e-03 & 0.373 & 0.484 & \textbf{0.931} & 0.9434 & 0.995 & 0.9871 \\
case300 & 7.746e-03 & 1.149e-02 & 1.865 & 2.691 & \textbf{0.935} & 0.9388 & 1.002 & 0.9897 \\
case9 & 2.642e-03 & 3.833e-03 & 0.326 & 0.461 & \textbf{0.968} & 0.9852 & 0.978 & 0.9908 \\
case14 & 1.759e-03 & 2.499e-03 & 0.230 & 0.302 & \textbf{0.973} & 0.9961 & 0.997 & 0.9975 \\
case30 & 1.812e-03 & 2.410e-03 & 0.215 & 0.286 & \textbf{0.980} & 0.9919 & 0.949 & 0.9631 \\
case5 & 1.906e-03 & 2.608e-03 & 0.113 & 0.148 & \textbf{0.982} & 0.9892 & 0.997 & 0.9952 \\
case118 & 2.324e-03 & 2.941e-03 & 0.490 & 0.633 & \textbf{0.983} & 0.9918 & 0.992 & 0.9949 \\
case4gs & 1.474e-03 & 1.695e-03 & 0.073 & 0.095 & \textbf{0.985} & 0.9990 & 0.999 & 0.9979 \\
case24\_ieee\_rts & 1.547e-03 & 2.055e-03 & 0.402 & 0.525 & \textbf{0.993} & 0.9965 & 1.002 & 0.9986 \\
case57 & 6.596e-03 & 8.646e-03 & 0.331 & 0.457 & \textbf{0.993} & 0.9945 & 0.999 & 0.9951 \\
SimBench & 8.297e-04 & 1.048e-03 & 0.022 & 0.034 & \textbf{0.998} & 0.9986 & 1.000 & 1.0000 \\
case\_ieee30 & 3.268e-03 & 4.530e-03 & 0.270 & 0.356 & \textbf{1.004} & 0.9844 & 0.999 & 0.9964 \\
case6ww & 1.099e-03 & 1.559e-03 & 0.199 & 0.269 & \textbf{1.008} & 0.9994 & 1.003 & 0.9855 \\
GBreducednetwork & 1.691e-02 & 2.289e-02 & 1.016 & 1.271 & \textbf{2.428} & 0.2825 & 0.994 & 0.9973 \\
\bottomrule
\end{tabular}
\caption{Re-scored from the saved checkpoints, sorted by $|V|$ regression slope
(worst collapse first). Inference only: the same parquet and the same test
split, so the RMSE reproduces the training run's --- verified on case118, where
this pass returns 0.011436366 against a logged 1.1436e-02. 27 of 31 grids;
the four largest were still re-scoring.}
\end{table}

\paragraph{Half the corpus is not tracking the voltage profile.} Of 27 grids,
13 have a $|V|$ slope below 0.9 and \textbf{6 are below 0.5} --- case3120sp
(0.208), LVN\_heo1 (0.213), case1354pegase (0.326), case2869pegase (0.330),
case89pegase (0.381), case1888rte (0.443). On these the model reproduces only
20--45\% of the reference voltage spread: it is regressing toward the mean
voltage and getting a respectable MAE for it. LVN\_heo1 is the sharpest case,
with a predicted standard deviation of 1.39e-2 pu against the reference's
3.11e-2.

\textbf{The RMSE table gives no hint of this.} LVN\_heo1 sits 16th of 30 on
RMSE, comfortably mid-pack, while explaining 23\% of the reference variance.

\paragraph{This corrects the earlier reading.} Section~5 concluded that ``error
lives in the angle, not the magnitude'', on the grounds that $\theta$ error
spans a far wider range than $|V|$ error. The regression says the opposite about
\emph{tracking}: the angle channel is healthy nearly everywhere
($\sin\theta$ slope median 0.992, range 0.776--1.003), while the magnitude
channel is what collapses. Both statements are true and they are about different
things --- the angle carries more absolute error but follows the reference,
whereas the magnitude carries less absolute error while not following it. The
magnitude is the weaker channel, and RMSE alone hid that.

\paragraph{One grid fails in the opposite direction.} GBreducednetwork has a
$|V|$ slope of \textbf{2.428} with $R^2 = 0.283$: it produces roughly two and a
half times the reference spread, and that excess is largely uncorrelated with
the reference. Amplitude inflation with low explained variance is a different
pathology from collapse, and only the slope/$R^2$ pair distinguishes them --- a
variance ratio alone would call this grid healthy.


\section{The released model, and a defect the comparison exposed}
\label{sec:graphkit}

The same 31 grids were then trained with IBM's released
\texttt{gridfm\_graphkit} model (\texttt{GNS\_heterogeneous}), which adds
per-layer physics feedback and a \texttt{PhysicsDecoderPF}. Neither
implementation has published weights, so both are scratch-only; this is a
comparison of architectures, not of pretraining.

\paragraph{They are not solving the same problem.} In AC power flow $|V|$ is an
\emph{input} at PV and slack buses and the angle is an input at the slack; only
the rest is unknown. The released model masks accordingly --- it copies those
values through and predicts a correction only where the quantity is free. The
local mirror does not: its \texttt{forward} never receives the mask its own
feature builder computes, and it adds a correction at every bus.

Measured on case118, where 54\% of buses are PQ:

\begin{table}[H]
\centering
\small
\begin{tabular}{lr}
\toprule
mirror, at the 46\% of buses where $|V|$ is given & \\
\midrule
mean $|V|$ moved away from the given value & 2.253e-03 pu \\
mean $|V|$ error against the solver        & 2.253e-03 pu \\
max                                        & 1.648e-02 pu \\
angle error at the slack bus               & 0.150 deg \\
\bottomrule
\end{tabular}
\caption{The first two rows are identical to every printed digit, and that is
the point: at PV and slack the given value \emph{is} the solution, so whatever
the model moves is exactly what it gets wrong. It is handed the answer and walks
away from it. The slack angle is the reference and should be exact.}
\end{table}

A \texttt{-{}-pin\_known} flag now holds these quantities to their inputs; a
verification run confirms the shift becomes exactly 0.000e+00 at PV/slack while
the PQ buses keep their full freedom. It is off by default so the campaign above
stays reproducible, and a separate pinned campaign is training.

\paragraph{Consequence for the comparison.} Pooled $|V|$ error is not
like-for-like: the released model is exact by construction at the buses it pins,
which on a grid like case5 is 80\% of them. The PQ-restricted columns are the
honest comparison.

\begin{table}[H]
\centering
\small
\begin{tabular}{lrrrrrrr}
\toprule
 & & \multicolumn{2}{c}{MAE $|V|$, all buses} & \multicolumn{2}{c}{slope $|V|$, all buses} & \multicolumn{2}{c}{graphkit, PQ only} \\
\cmidrule(lr){3-4}\cmidrule(lr){5-6}\cmidrule(lr){7-8}
Grid & PQ & mirror & graphkit & mirror & graphkit & MAE & slope \\
\midrule
SimBench & 97\% & 8.297e-04 & 1.279e-03 & 0.998 & 0.992 & 1.321e-03 & 0.992 \\
case6ww & 50\% & 1.099e-03 & 2.610e-04 & 1.008 & 1.001 & 5.219e-04 & 0.998 \\
case4gs & 50\% & 1.474e-03 & 1.731e-04 & 0.985 & 1.005 & 3.462e-04 & 1.003 \\
case14 & 64\% & 1.759e-03 & 5.233e-04 & 0.973 & 0.999 & 8.141e-04 & 0.993 \\
case5 & 20\% & 1.906e-03 & 2.108e-05 & 0.982 & 1.000 & 1.054e-04 & 0.999 \\
case118 & 54\% & 2.324e-03 & 6.041e-04 & 0.983 & 0.991 & 1.114e-03 & 0.979 \\
case9 & 67\% & 2.642e-03 & 4.954e-04 & 0.968 & 0.991 & 7.431e-04 & 0.986 \\
case\_ieee30 & 80\% & 3.268e-03 & 2.911e-03 & 1.004 & 0.970 & 3.639e-03 & 0.934 \\
case\_illinois200 & 81\% & 4.745e-03 & 5.466e-03 & 0.931 & 0.922 & 6.748e-03 & 0.860 \\
case39 & 74\% & 5.107e-03 & 5.638e-03 & 0.889 & 0.852 & 7.583e-03 & 0.805 \\
case57 & 88\% & 6.596e-03 & 3.798e-03 & 0.993 & 0.981 & 4.329e-03 & 0.974 \\
case33bw & 97\% & 7.655e-03 & 5.970e-03 & 0.872 & 0.891 & 6.157e-03 & 0.882 \\
case300 & 77\% & 7.746e-03 & 1.525e-02 & 0.935 & 0.809 & 1.981e-02 & 0.712 \\
iceland & 81\% & 1.316e-02 & 1.063e-02 & 0.897 & 0.856 & 1.304e-02 & 0.838 \\
GBnetwork & 83\% & 1.572e-02 & 2.363e-02 & 0.739 & 0.574 & 2.846e-02 & 0.549 \\
case1354pegase & 81\% & 1.887e-02 & 1.390e-02 & 0.326 & 0.472 & 1.720e-02 & 0.352 \\
case2848rte & 87\% & 2.012e-02 & 2.403e-02 & 0.621 & 0.442 & 2.761e-02 & 0.337 \\
LVN\_heo1 & 99\% & 2.052e-02 & 1.779e-02 & 0.213 & 0.503 & 1.791e-02 & 0.494 \\
case1888rte & 86\% & 2.443e-02 & 2.136e-02 & 0.443 & 0.461 & 2.495e-02 & 0.228 \\
\bottomrule
\end{tabular}
\caption{19 grids re-scored under both implementations, sorted by mirror MAE.
The mirror has no PQ-restricted columns: its rescore predates the metric, and
re-running it is pending.}
\end{table}

\paragraph{The released model wins where the pinning helps most.} It is ahead on
13 of 19 grids. The margin is largest exactly where the pinned fraction is
largest: case5 (20\% PQ) by $90\times$, case4gs (50\%) by $8.5\times$,
case9 (67\%) by $5.3\times$ and case118 (54\%) by $3.8\times$. On grids that are almost
entirely PQ --- SimBench at 97\%, LVN\_heo1 at 99\% --- the two are within a
factor of 1.5 either way. That pattern is what one would expect if a large part
of the measured advantage is the masking rather than the architecture.

\paragraph{Neither architecture escapes the collapse.} Counting $|V|$ slope
below 0.9: 8 of 19 for the mirror, 9 of 19 for the released model. Both
regress toward the mean voltage on the same hard grids --- case1888rte
(0.443 / 0.461), case2848rte (0.621 / 0.442), GBnetwork (0.739 / 0.574). The
physics decoder does not fix it, and on GBnetwork and case2848rte the released
model is \emph{worse}. LVN\_heo1 is the one clear reversal, where it more than
doubles the tracked spread (0.213 $\rightarrow$ 0.503).

Restricted to PQ buses the released model's slope falls further still: 0.228 on
case1888rte, 0.337 on case2848rte and 0.352 on case1354pegase. The pinned
buses, trivially perfect, had been propping the pooled figure up.

\section{Caveats}

\paragraph{case145 is not comparable.} It is the worst grid on voltage
(RMSE $2.06\times10^{-1}$, $\theta = 11.7^\circ$) and fourth worst on residual,
and it is also the only grid with 7885 scenarios instead of 36{,}000 --- 22\% of
the data everyone else got. A $\pm20\%$ regeneration is outstanding; the grid
should be re-run before its number is used for anything.

\paragraph{Scenario counts are not uniform.} Besides case145: case9241pegase
18{,}280, case6515rte 20{,}135, case6495rte 23{,}332, case6470rte 25{,}551,
case300 27{,}994, case1354pegase 32{,}627, GBnetwork 34{,}424, iceland 35{,}605.
Cross-grid comparisons carry this confound.

\paragraph{Coverage.} The mirror campaign is complete (31/31) and 30 of 31 are
re-scored. The released-model campaign is 30/31 trained with 19 re-scored, so
Section~\ref{sec:graphkit} compares the 19 grids both have. A third campaign,
the mirror with \texttt{-{}-pin\_known}, is training.

\paragraph{No pretrained baseline exists.} \texttt{gridfm\_graphkit} ships no
weights, so GridFM is scratch-only under either implementation. Unlike LUMINA
and GridSFM in the OPF study, there is no ``pretrained'' variant to compare
against, and none of these runs should be described as fine-tuned.

\end{document}
